
\begin{frame}[fragile]
  \frametitle{Evolución del Recaudo de los Principales Impuestos (\% del PIB)}

%  \scriptsize
  \begin{table}
  \centering
  \begin{tabular}{|l|c|c|c|c|}
  \hline
  \textbf{Año} & \textbf{Renta*} & \textbf{IVA} & \textbf{Resto} & \textbf{Total} \\
  \hline
  2012 & 6,6 & 5,5 & 2,3 & \textbf{14,3} \\
  2013 & 6,7 & 4,9 & 2,6 & \textbf{14,1} \\
  2014 & 6,6 & 5,1 & 2,5 & \textbf{14,2} \\
  2015 & 6,6 & 5,2 & 2,7 & \textbf{14,5} \\
  2016 & 6,8 & 4,3 & 2,4 & \textbf{13,6} \\
  2017 & 6,8 & 4,9 & 2,0 & \textbf{13,8} \\
  2018 & 7,1 & 5,0 & 1,5 & \textbf{13,7} \\
  2019 & 7,0 & 5,2 & 1,8 & \textbf{14,0} \\
  2020 & 6,8 & 4,8 & 1,5 & \textbf{13,1} \\
  2021 & 6,8 & 5,2 & 1,6 & \textbf{13,6} \\
  2022 & 7,2 & 5,5 & 1,7 & \textbf{14,4} \\
  2023 & 9,6 & 5,3 & 1,8 & \textbf{16,7} \\
  2024 & 7,5 & 5,0 & 1,9 & \textbf{14,4} \\
  \hline
  \end{tabular}
%  \vspace{0.1cm}

  \par\footnotesize{* Incluye recaudo del CREE entre 2013 y 2016.}
%  \par\footnotesize{\textit{Fuente: Balance General del GNC, Ministerio de Hacienda.}}
  \end{table}

\end{frame}
